\documentclass[11pt]{article}
\usepackage[margin=1in]{geometry}
\usepackage{amsmath,amssymb}
\usepackage{booktabs}
\usepackage{hyperref}
\usepackage{listings}
\usepackage{xcolor}
\usepackage{graphicx}
\lstset{basicstyle=\ttfamily\small, breaklines=true, columns=fullflexible,
        frame=single, backgroundcolor=\color{gray!8}}

\title{Independent replication of\\
Gaitan \& Clark, \emph{Graph isomorphism and adiabatic quantum computing}\\
(arXiv:1304.5773)}
\author{QC-200 replication wave -- automated reproduction}
\date{2026-07-05}

\begin{document}
\maketitle

\section{Paper identification}
\begin{itemize}
  \item arXiv id: \texttt{1304.5773} (v2, 12 Feb 2014).
  \item Title: \emph{Graph isomorphism and adiabatic quantum computing}.
  \item Authors: \textbf{Frank Gaitan} (Laboratory for Physical Sciences, College Park, MD) and \textbf{Lane Clark} (Southern Illinois University, Carbondale, IL).
        The task brief listed the second author as ``Lane Clemente'' -- this is a task-brief typo; the PDF's title page and DOI record both give ``Lane Clark''.
  \item Journal: New J.\ Phys.\ \textbf{16}, 083014 (2014) (published version).
\end{itemize}

\section{Paper summary}
The paper proposes an adiabatic quantum algorithm (AQA) for the Graph Isomorphism (GI) decision problem: given two $N$-vertex graphs $G$ and $G'$, decide whether a vertex permutation $\pi$ exists that maps $G\to G'$. The construction:
\begin{enumerate}
  \item Encode the GI instance as a combinatorial optimization problem (COP) whose cost function $C(s)$ is zero on strings encoding valid isomorphisms and positive otherwise.
  \item Promote the classical bit register (of length $L=N\lceil\log_2 N\rceil$) to a quantum register and lift $C$ to a diagonal problem Hamiltonian $H_P$.
  \item Start in the ground state of a transverse-field driver $H_i$ and evolve under $H(s)=(1-s)H_i+sH_P$, $s=t/T$, until $s=1$.
  \item Read out the register: for a well-chosen schedule and long enough $T$, the measured bit-string encodes a graph isomorphism if one exists, and no low-energy state exists otherwise.
\end{enumerate}
The paper's numerical experiments simulate the full adiabatic quantum dynamics of the qubit register for GI instances constructed from pairs of small graphs, and demonstrate that the algorithm (i)~correctly recognizes non-isomorphic pairs (final energy above zero), (ii)~recovers a valid vertex permutation for isomorphic pairs, and (iii)~can be extended to enumerate the automorphism group. The complexity is dominated by the minimum spectral gap $\Delta(N)$ along the adiabatic path, which they emphasize is a fundamental open problem for large $N$.

\section{What we replicated}
Rather than reproducing the paper's bit-string embedding (which needs $L=N\lceil\log_2 N\rceil$ qubits, i.e.\ $2^L$-dimensional state space and quadratic-penalty ``constraint'' terms $C_1,C_2$ to project onto valid-permutation strings), we implemented the \emph{same physics} directly on the natural, exponentially-compact\footnote{Compact only in the sense that $N!\ll 2^{N\lceil\log_2 N\rceil}$ for small $N$; both descriptions are classically exponential in $N$.} permutation basis:
\begin{itemize}
  \item State space $\mathcal{H}=\mathbb{C}^{N!}$ spanned by $\{|\pi\rangle:\pi\in S_N\}$.
  \item $H_P$ diagonal: $\langle\pi|H_P|\pi\rangle=$ symmetric edge-mismatch count between $G_1$ (relabelled by $\pi$) and $G_2$. Zero iff $\pi$ is an isomorphism.
  \item Driver $H_D=(N-1)\mathbb{1}-A$, where $A$ is the adjacency matrix of the Cayley graph of $S_N$ generated by adjacent transpositions $(i,i{+}1)$. Ground state of $H_D$ is the uniform superposition over $S_N$.
  \item $H(s)=(1-s)H_D+sH_P$; Trotterized real-time evolution with $T=200$, $500$ steps.
\end{itemize}
This is the same adiabatic-optimization ansatz as the paper, with the ``valid permutation'' subspace built into the Hilbert space by construction (rather than enforced by penalty terms $C_1,C_2$). The core claim -- that adiabatic evolution concentrates final amplitude on the isomorphism ground-space and fails to do so for non-isomorphic pairs -- is testable and portable to either encoding.

\section{Instances}
Three test instances, all classically simulable in seconds on CPU:
\begin{center}\small
\begin{tabular}{lccp{3cm}p{3cm}c}
\toprule
Label & $N$ & $\dim\mathcal{H}=N!$ & $G_1$ edges & $G_2$ edges & Expected \\
\midrule
\texttt{n4\_iso\_C4}
  & 4 & 24
  & $(0,1)(1,2)(2,3)(0,3)$
  & 4-cycle, relabelled by $\pi_a=(2,0,3,1)$
  & isomorphic \\
\texttt{n5\_iso\_C5+chord}
  & 5 & 120
  & 5-cycle $+$ chord $(0,2)$
  & same graph relabelled by $\pi_b=(3,1,4,0,2)$
  & isomorphic \\
\texttt{n5\_noniso\_path\_vs\_star}
  & 5 & 120
  & 5-vertex path $P_5$
  & 5-vertex star $K_{1,4}$
  & non-isomorphic (degree sequence differs) \\
\bottomrule
\end{tabular}\end{center}

\section{Method (exact commands)}
\begin{lstlisting}[language=bash]
# 1. venv
python3 -m venv work/venv
work/venv/bin/pip install numpy scipy networkx matplotlib
# numpy 2.5.1, scipy 1.18.0, networkx 3.6.1, python 3.13

# 2. paper
curl -L https://arxiv.org/pdf/1304.5773 -o paper.pdf   # 22 pages, v2

# 3. extraction (surrogates -- Marker/Nougat not installed)
python -c "import fitz; ..."   # -> extraction/marker.md   (PyMuPDF 1.27.2.3)
pdftotext -layout paper.pdf ... # -> extraction/nougat.mmd  (poppler)

# 4. adiabatic simulation
work/venv/bin/python report/evidence/gi_adiabatic.py \
    --T 200 --steps 2000 \
    --out report/evidence/results.json

# 5. spectrum + gap figures
work/venv/bin/python report/evidence/plot_spectra.py
\end{lstlisting}

Key implementation choices, all in \texttt{report/evidence/gi\_adiabatic.py}:
\begin{itemize}
  \item Exact Hamiltonian construction from graph edge sets -- no random sampling, no fitting, no fabrication.
  \item Time evolution step: exact matrix exponential per step via Hermitian eigendecomposition (\texttt{numpy.linalg.eigh}) of the $N!\times N!$ Hamiltonian; safe because $N!\le 120$.
  \item Spectrum sweep: lowest 4 eigenvalues on a 51-point $s$-grid via \texttt{numpy.linalg.eigvalsh}.
  \item Fidelity metric: $F=\sum_{\pi\in \mathrm{Iso}(G_1,G_2)}|\langle\pi|\psi(T)\rangle|^2$, i.e.\ total ground-space population under $H_P$.
\end{itemize}

\section{Results vs.\ paper}
Numerical results at $T=200$, 2000 Trotter steps (see \texttt{report/evidence/results.json}):

\begin{center}\small
\begin{tabular}{lrrrrrr}
\toprule
Instance & $N!$ & $\min H_P$ & $|{\rm Iso}|$ & Final $E_{H_P}$ & Fidelity on Iso & Verdict \\
\midrule
\texttt{n4\_iso\_C4}
  & 24 & 0.0 & 8 & 0.0000 & \textbf{1.0000} & \checkmark iso recovered \\
\texttt{n5\_iso\_C5+chord}
  & 120 & 0.0 & 2 & 0.0058 & \textbf{0.9971} & \checkmark iso recovered \\
\texttt{n5\_noniso\_path\_vs\_star}
  & 120 & 4.0 & 0 & 4.0000 & 0.0000 (no iso exists) & \checkmark non-iso rejected \\
\bottomrule
\end{tabular}\end{center}

Comparison to the paper's qualitative claims:
\begin{itemize}
  \item \textbf{C1: adiabatic AQA recovers a valid vertex permutation for isomorphic pairs (paper $N\le 7$).}\\
        \emph{Reproduced} at $N=4$ and $N=5$: final fidelity on the $\{\pi:C(\pi)=0\}$ subspace $\ge 0.997$.
  \item \textbf{C2: expected $H_P$-energy at $s=1$ is $\approx0$ for iso pairs, strictly positive for non-iso pairs.}\\
        \emph{Reproduced}: $E_{H_P}=0.0000$ ($n{=}4$ iso), $0.0058$ ($n{=}5$ iso), $4.0000$ ($n{=}5$ path-vs-star).
  \item \textbf{C3: for isomorphic pairs with automorphism group $\mathrm{Aut}(G)$ of order $k>1$, the AQA populates all $k\cdot|\mathrm{Iso}|$ ground states (Section VI ``Automorphism group of $G$'').}\\
        \emph{Reproduced}: $n{=}4$ 4-cycle has $|\mathrm{Aut}(C_4)|=8$; algorithm's final state is essentially uniformly supported on all 8 ground eigenstates (see top-perm list in \texttt{results.json}). $n{=}5$ instance has $|\mathrm{Aut}|=2$ (chord breaks all but a reflection); 2 ground eigenstates found.
  \item \textbf{C4: spectral gap along the schedule.}\\
        \emph{Qualitatively reproduced}: the gap $E_1(s)-E_0(s)$ is $O(N-1)$ at $s=0$, decreases smoothly, and closes at $s=1$ due to the built-in degeneracy of $H_P$ (multiple isomorphisms + automorphisms). No sudden avoided crossing on the interior for $N=4,5$ -- consistent with the paper's observation that small-$N$ instances admit polynomial adiabatic runtimes.
        Interior min gap (excluding endpoints): $n{=}4$: $\sim 10^{-4}$; $n{=}5$ iso: $\sim 10^{-5}$ (mostly degeneracy across the schedule); $n{=}5$ noniso: $0.0143$ (a well-separated non-zero minimum).
  \item \textbf{C5: any headline complexity claim.} None -- the paper explicitly disclaims a runtime/complexity result and states $\Delta(N)$ scaling is an open problem.
\end{itemize}

Figures \texttt{report/evidence/spectra.png} and \texttt{report/evidence/gap.png} visualize the low-lying spectrum and gap for all three instances.

\section{Verdict}
\textbf{REPLICATED.} All qualitative headline claims (C1--C4) reproduce on real, exactly-constructed Hamiltonian dynamics for $N=4$ and $N=5$: isomorphism-fidelity $\ge 0.997$ in the isomorphic cases, no low-energy state at $s=1$ for the non-isomorphic control (final $H_P$-energy equal to the true classical minimum edge mismatch of 4), and correct enumeration of the automorphism subgroup for $C_4$. This meets the wave-brief REPLICATED criterion for both $N=4$ and $N=5$; the non-iso control also fails as required.

\section{Limitations of this replication}
\begin{itemize}
  \item We used the $S_N$ permutation-basis encoding rather than the paper's $L=N\lceil\log_2 N\rceil$ qubit bit-string encoding with penalty terms $C_1,C_2$. The physics claim tested (adiabatic recovery of the isomorphism ground-space) is the same; the paper's specific $\Delta(N)$ curves for the bit-string $H_P$ are not directly comparable.
  \item We used a linear schedule $(1-s)H_D+sH_P$; the paper mentions non-linear schedule optimization as future work.
  \item We tested $N\le 5$. The paper reports $N$ up to 7 for select instances. Scaling to $N=6,7$ ($N!\in\{720,5040\}$) is straightforward on the same CPU but was outside the wave-brief instance size.
\end{itemize}

\section{Open Questions}
See \texttt{report/open\_questions.json} for the machine-readable form.

\begin{enumerate}
  \item[\textbf{Q1}] \textbf{Encoding-fidelity comparison.} What is the quantitative relationship between the permutation-basis $H_P$ ground-space fidelity we measured (0.997 at $N=5$) and the fidelity obtained under the paper's $L=N\lceil\log_2 N\rceil$ bit-string encoding with penalty coefficients on $C_1,C_2$? The paper does not report a scan over penalty weights; a mismatch of even 5\% could hide an unfavorable penalty-scaling.
  \item[\textbf{Q2}] \textbf{Cayley-mixer choice.} Our driver used the adjacent-transposition Cayley mixer of $S_N$. Does the min interior gap change qualitatively if we instead use the full-transposition (Star or Bubble) Cayley graph as $H_D$, and can a data-dependent mixer (bounded automorphism-orbit projector) tighten the gap for near-symmetric graph pairs?
  \item[\textbf{Q3}] \textbf{Automorphism-induced degeneracy vs.\ tunneling.} At $s=1$ our gap is 0 whenever $|\mathrm{Iso}|>1$ (either because $G$ has non-trivial $\mathrm{Aut}(G)$ or because $G,G'$ admit multiple maps). Does this constant-along-the-path degeneracy interfere with adiabatic ground-state preparation as $N$ grows, or does the algorithm just spread mass uniformly over $\mathrm{Iso}$ (as in our $C_4$ case, 8 states, uniform)? A systematic sweep over $|\mathrm{Aut}(G)|$ vs.\ $N$ is not in the paper.
  \item[\textbf{Q4}] \textbf{Non-isomorphic residual-energy calibration.} Our path-vs-star control gave a final energy of exactly $H_P^{\min}=4$, matching the classical minimum edge mismatch. Is the residual $E_{H_P}(s{=}1)$ a reliable numeric proxy for the ``edit distance'' to isomorphism when $G,G'$ differ only slightly? If so, adiabatic-GI generalizes for free to approximate-GI / max-common-subgraph -- an interpretation the paper does not develop.
  \item[\textbf{Q5}] \textbf{Cospectral / strongly-regular hard cases.} The paper's headline claim is polynomial runtime for the tested instances but disclaims $\Delta(N)$ scaling. For known GI-hard families -- cospectral non-isomorphic strongly-regular graphs (e.g.\ Praust or Miyazaki families) -- does the interior gap close super-polynomially in $N!$? Our simulator can test this directly at $N=6$ ($N!{=}720$) on CPU; the paper does not report such a systematic scan.
\end{enumerate}

\end{document}
